% Abstract (English)
% Default language ditangani di preamble. Tidak perlu switch di sini.
\begin{center}
{\bfseries\large ABSTRACT}
\end{center}

\vspace{0.5cm}

\begin{spacing}{1}

\noindent {\bfseries SOFTWARE ARCHITECTURE DEVELOPMENT OF A COMPANIONSHIP CHATBOT WITH LONG-TERM MEMORY BASED ON KNOWLEDGE GRAPH FOR INDONESIAN STUDENTS}

\vspace{0.3cm}

\noindent By \\
\noindent Mohammad Nugraha Eka Prawira \\
\noindent NIM: 13522001

\vspace{0.5cm}

\noindent \textit{Placeholder --- abstract left empty for now since the current focus is the seminar hasil (Chapters I--IV and related appendices). The previous content is preserved in the \textbackslash iffalse block below and will be redone ahead of the final defense.}

\vspace{0.5cm}

\noindent {\bfseries Keywords:} companionship chatbot, long-term memory, knowledge graph, guardrail system, Cognitive Behavioral Therapy, Indonesian university students
\end{spacing}

\iffalse
% =====================================================================
% Original abstract content (kept for reuse ahead of the final defense)
% =====================================================================
Indonesian university students face considerable academic and emotional pressure, while access to professional support remains limited, leaving room for accessible, nonclinical forms of support. Existing companionship chatbots typically excel at only one dimension: some are strong at reflective techniques, others at conveying a sense of companionship, but few combine empathetic conversation, mood screening, and long-term memory that consistently links a user's experiences across sessions. This thesis develops AXIS, a nonclinical companionship chatbot for Indonesian university students. The three research questions concern safe and empathetic conversation that fits Indonesian students' language use, the integration of daily mood check-ins and PHQ-9 screening into the companionship flow, and knowledge-graph-based long-term memory that can write, link, update, and re-surface user experiences relevantly over time. The solution is realized as a decision-graph-based agentic pipeline with a four-layer guardrail, a CBT dialogue policy, PHQ-9 modeled as a conversational state machine, and a hybrid memory architecture combining a relational knowledge graph with semantic search. Evaluation proceeds along three tracks with differing levels of maturity: functional validation of critical components, an illustrative case study comparing AXIS against a baseline chatbot built on pure semantic retrieval, and a design for a limited real-user study. Functional validation and case-study results provide early indication that the proposed architecture behaves as designed, with methodological limitations stated explicitly, particularly regarding the small case-study sample size and the real-user validation that had not yet been conducted at the time this report was written.
\fi

% Default language ditangani di preamble.
